%% Introduction %%

Pour terminer sur cette accumulation de facteurs résultant en cette absence de visibilité sur le \gls{td} dans les bases de données actuelles de l'INA. Il est important de souligner la quantité massive de données que conserve l'INA. Nous avons déjà évoqué à plusieurs reprises cette quantité massive de données. Afin de comprendre en quoi cela affecte l'accessibilité aux données, nous allons nous attarder dessus un temps. 

 %% Sources des données %%

C'est un fait, l'INA conservent une quantité très importantes de données, elles sont le résultat d'une collecte provenant de sources multiples et très productrice de données. Nous en avons cité les principales avec la captation au titre du dépôt légal, elle enregistre 24 heures sur 24 120 chaînes télévisées, ce qui ajoute aux collections numériques 1 millions d'heures par an.  A cela s'ajoute les fonds de l'ORTF les Actualités françaises, les fonds issus de divers dons, le dépôt légal du web etc... 

%% PSN %%

Mais le premier et principal déclencheur de cette quantité de données est le PSN, Plan de Sauvegarde Numérique déclenché en 1999. Nous sommes dans une période où la numérisation se répand dans les institutions patrimoniales. Dans beaucoup, il est nécessaire de sélectionner les documents à numériser. C'est le cas à la BNF qui choisit les documents de ses collections à numériser en fonction de critères juridiques, techniques et intellectuels\footcite[][155]{bermès2024}. L'INA ne fait pas face aux même problématiques concernant la sélection des documents numérisés. Au contraire, les supports physiques qu'ils conservent atteignent leur limites. Ils se dégradent rapidement, ce qui ne leur permet pas d'être exploités d'aucune façon, certaines émissions n'étant conservés qu'en un seul exemplaire. Ce à quoi nous pouvons ajouter les équipements de lecture pouvant manquer ou amenés à disparaître dans un futur proche pour certains supports. S'ajoutaient à la difficulté d'exploitation de ces supports vieillissants la difficulté et le coût pour faire parvenir ces documents pour leur exploitation\footcite[28]{varra2023}. Ainsi, des années durant, au cours d'une campagne qui devait durer 15 ans mais qui s'étendit bien au delà, l'INA numérise ses collections. C'est un pari gagnant, la communication est plus aisée, les documents échappent à la disparition totale mais en contrepartie l'INA rentre dans ce que l'on appelle le \textit{Big Data}. 

%% Le Big Data %%

Le Big data, que l'on peut traduire en \enquote{données massives} en français identifie les bases de données contenant un nombre très important de données, souvent hétérogènes nécessitant de nouvelles technologies accordée à cette quantité de données pour en faire ressortir des informations. Ce terme est devenu très populaire dans le milieu des bibliothèque grâce à ce mouvement de numérisation dont le projet \enquote{Google Books}\footcite[]{demauro2016}, numérisant une quantité astronomique de livres, en est un exemple. Tout comme le sont la campagne de numérisation des collections de la BNF et de l'INA que nous venons d'évoquer. Ces deux campagnes correspondent parfaitement à cette définition car elles contribuent à rassembler, par la numérisation, des quantités massives de données, hétérogènes à cause de leurs origines (captation, numérisation), de leur support d'origine (physique : livres, CD, carte, estampe ou nativement numérique) ou même de l'objectif de leur collecte (conservation, communication). C'est pourquoi les méthodes de catalogage et de \gls{td} telles que pensées par Paul Otlet et Suzanne Briet ne correspondent plus aux nouvelles collections constituées par les institutions patrimoniales. Elles nécessitent de nouveaux outils et de nouvelles méthodes, autant pour être stockées que pour être exploitées. En conséquence, l'échelle du \gls{td} change drastiquement, leurs usages également. L'information peut circuler plus rapidement alors les pratiques métiers doivent s'adapter\footcite[][145]{caporiondo-prost2026}. Le \gls{td} doit être effectué plus rapidement qu'il ne l'avait jamais été et pourtant, les fonds doivent rester cohérents malgré cette hétérogénéité. 

%% Conclusion %%

Si les collections des institutions patrimoniales n'avaient que les incohérences que nous avons citées précédemment mais que la quantité de données était plus réduite, l'accessibilité aux collections serait facilitée. Mais dans le cas du \enquote{Big Data}, ces erreurs et incohérences créent ces données hétérogènes en une quantité plus difficilement gérable. Cette quantité massive de données peut provoquer un bruit important lors de la recherche, surtout dans notre cas où la cartographie du traitement de ces données n'est pas clair. Les documentalistes ont l'habitude d'être confrontées à ces données mais ils ont conscience que cela peut être un obstacle à la recherche et mettent en garde les personnes nécessitant de faire des recherches dans les bases, dont les chercheurs\footnote{(Voir entretien d'une documentaliste de l'INA sur l'impact de la \gls{pd} sur les données de description et de Pascal Lebeurier sur l'expertise du documentaliste ~\ref{ann:transcriptions},\nameref{ann:transcriptions}), p.~\pageref{ann:transcriptions}.}. A cette occasion, l'expertise du documentaliste devient essentielle. Il sait autant quels pièges éviter lors de la recherche que quelles informations rechercher et dans quels champs. Dans son entretien, Pascal Lebeurier souligne que sans méthode, il est aisé de se noyer dans cette masse de données et Arthur Hénaff souligne l'importance que prend la médiation de données à une époque où les quantité de données sont si démesurées, d'autant plus maintenant que les résultats des traitements IA se sont ajoutés.